(1) 横向力平衡为 \(\mu y_{tt}=\partial_x[T(x)y_x]\)，设 \(y=Y(x)e^{i\omega t}\)，得 \(-(TY')'=\mu\omega^2Y\) 与 \(Y(0)=Y(L)=0\)。乘以 \(Y\) 积分并分部积分后，任一满足边界的试探函数对应 Rayleigh 商 \(\mathcal R[Y]=[\int_0^L T(x)Y'^2dx]/[\int_0^L\mu Y^2dx]\)，正定且具有角频率平方的量纲。
(2) 对 \(Y_{\rm tr}=\sin(\pi x/L)\)，分母为 \(\mu L/2\)；分子为 \(T_0\pi^2/L^2\int_0^L(1+\alpha x/L)\cos^2(\pi x/L)dx=T_0\pi^2(1+\alpha/2)/(2L)\)，利用 \(\int_0^L x\cos^2(\pi x/L)dx=L^2/4\)。故 \(\omega_{\rm tr}^2=(T_0/\mu)(\pi^2/L^2)(1+\alpha/2)\)。
(3) 自伴 Sturm–Liouville 算子的基频平方满足 \(\omega_1^2=\min_{Y\ne0}\mathcal R[Y]\)。由于所选试探函数只是所有允许函数之一，必有 \(\omega_1^2\le\omega_{\rm tr}^2\)，这是上界而绝非下界；若 \(\alpha=0\)，该正弦函数就是精确基模，估计取等号。
(4) 因 \(T_{\min}=T_0\min(1,1+\alpha)\)、\(T_{\max}=T_0\max(1,1+\alpha)\)，对全部试探函数均有 \(T_{\min}\int Y'^2\le\int TY'^2\le T_{\max}\int Y'^2\)。结合固定端 Poincaré 不等式及常张力正弦试探，得 \((T_{\min}/\mu)(\pi/L)^2\le\omega_1^2\le(T_{\max}/\mu)(\pi/L)^2\)，更紧的上界是第(2)问的 \(\omega_{\rm tr}^2\)。当 \(\alpha\to-1^+\) 时一端张力接近零，真实基模会改变形状，正弦试探不保证继续紧。
